\documentclass[11pt]{article}
\usepackage[margin=1in]{geometry}
\usepackage{amsmath,amssymb}
\usepackage{enumitem}

\title{Summary: Unitizing Krippendorff $\alpha_u$ (character continuum) in \texttt{iaa\_scores}}
\date{}

\begin{document}
\maketitle
\vspace{-1.2em}

\noindent This note summarizes how this repository computes Krippendorff's unitizing $\alpha_u$
on a character-offset continuum. It focuses on the implemented procedure in
\texttt{iaa\_scores/alpha.py} (\texttt{compute\_alpha\_u\_per\_doc}, \texttt{compute\_alpha\_u\_overall}).

\section*{Extracted data and notation}
\begin{itemize}[leftmargin=*,nosep]
  \item Let $d$ be a document (identified by \texttt{ref\_id}) and $a\in\mathcal{A}_d$ annotators for $d$.
  \item Let $\mathcal{C}$ be the set of label categories observed in the corpus.
  \item For an annotator $a$ and category $c$, the loader extracts offset intervals
  \[
    I_{a,c}(d)=\{(s,e)\}
  \]
  from Label Studio \texttt{result} items of type \texttt{"labels"} (requiring valid integer offsets and $s<e$;
  see \texttt{iaa\_scores/loaders.py}).
  \item When available, the continuum length is taken as the document text length
  $L_d = |\texttt{task.data.case\_content}|$.
\end{itemize}

\section*{Continuum length and per-category segmentation}
For each document, the continuum length used by $\alpha_u$ is
\[
L \;=\; \max\Big(L_d,\ \max\{e : (s,e)\in I_{a,c'}(d)\text{ for some }a\in\mathcal{A}_d,\ c'\in\mathcal{C}\}\Big),
\]
so $L$ is at least the largest annotated end offset, but can extend to the end of the document text.

Fix a category $c$. For each annotator $a$, the interval list $I_{a,c}(d)$ is merged into disjoint
intervals (union of overlapping or touching spans; \texttt{iaa\_scores/alpha.py:\_merge\_intervals}).
Define the boundary set
\[
B \;=\; \{0,L\}\cup\{s,e : (s,e)\in I_{a,c}(d)\text{ for some }a\}
\quad\text{(\texttt{iaa\_scores/alpha.py:\_boundaries}).}
\]
Let $0=b_0<b_1<\cdots<b_K=L$ be the sorted boundaries; define segments $t_k=[b_k,b_{k+1})$ with length
$\ell_k=b_{k+1}-b_k$. For each annotator $a$, define a binary coverage indicator on segments:
\[
x_a(k)\;=\;\mathbf{1}\big[\exists(s,e)\in I_{a,c}(d)\text{ such that } s\le b_k < e\big]
\quad\text{(\texttt{iaa\_scores/alpha.py:\_covers}).}
\]
Thus, for a fixed category $c$, each character-offset segment is treated as either
``inside a $c$ span'' ($1$) or ``not-$c$'' ($0$). Repeated text is irrelevant: units are defined
by offset position, so the same phrase appearing in two different places is two different units.

\section*{Coincidence matrix (length-weighted)}
For each unordered annotator pair $(a_i,a_j)$ with $i<j$, the code accumulates segment lengths into
$O\in\mathbb{R}^{2\times 2}$ as:
\[
O_{x_{a_i}(k),\,x_{a_j}(k)} \;\leftarrow\; O_{x_{a_i}(k),\,x_{a_j}(k)} + \ell_k,
\qquad
O_{x_{a_j}(k),\,x_{a_i}(k)} \;\leftarrow\; O_{x_{a_j}(k),\,x_{a_i}(k)} + \ell_k,
\]
but only if \texttt{include\_background} is true or $x_{a_i}(k)=1$ or $x_{a_j}(k)=1$
(\texttt{iaa\_scores/alpha.py:\_coincidence\_lengths\_binary}).
Intuitively:
\begin{itemize}[leftmargin=*,nosep]
  \item $O_{1,1}$ is the total length where both annotators label $c$,
  \item $O_{1,0}$ and $O_{0,1}$ are the total length where exactly one annotator labels $c$,
  \item $O_{0,0}$ is the total length where neither annotator labels $c$ (only counted when \texttt{include\_background=true}).
\end{itemize}

\section*{From coincidence matrix to $\alpha_u$ (repo formula)}
Given $O$, define $n=\sum_{p,q\in\{0,1\}}O_{pq}$ and row marginals $m_p=\sum_{q}O_{pq}$
(\texttt{iaa\_scores/alpha.py:\_alpha\_from\_coincidence}). The code computes an ``expected'' matrix $E$ via
\[
E_{pp}=\frac{m_p(m_p-1)}{\max(n-1,1)},\qquad
E_{pq}=\frac{m_pm_q}{\max(n-1,1)} \ (p\ne q)
\quad\text{(\texttt{iaa\_scores/alpha.py:\_alpha\_from\_coincidence}).}
\]
With nominal binary distance $\delta(p,q)=\mathbf{1}[p\ne q]$, it sets
\[
D_o=\sum_{p\ne q} O_{pq},\qquad D_e=\sum_{p\ne q} E_{pq}
\quad\text{(\texttt{iaa\_scores/alpha.py:\_alpha\_from\_coincidence}),}
\]
and returns $\alpha_u=1-D_o/D_e$, but returns $1.0$ if $n=0$ or $D_e=0$
(\texttt{iaa\_scores/alpha.py:\_alpha\_from\_coincidence}).

\section*{\texttt{include\_background} (recommended)}
For per-category binary unitizing, ``background'' corresponds to the value $0$ (not-$c$), which includes
both unlabeled text and text labeled with other categories.
Setting \texttt{include\_background=true} counts $(0,0)$ segment agreements and therefore measures agreement
about where category $c$ is \emph{absent} as well as present. This is preferred when annotators are expected
to agree not only on the positive spans, but also on leaving other regions unlabeled (or labeled differently).

\end{document}
